% RETAINED LEGACY SOURCE: not included by main.tex after the structural audit.
% Active replacement: sections/interpretation.tex and appendices/quotation_dossier.tex
\section{Interpreting Genetic Estimates: Goldberger and Jencks}
\label{sec:interpretation}

The preceding criticisms concern the identification and estimation of
Clark's genetic model. A separate question remains even if those estimates
are accepted: what do they establish about social mechanisms, opportunity,
and policy? Two classical critiques distinguish the issues particularly
clearly. Clark recognizes important parts of both, but does not consistently
carry those qualifications through to his stronger conclusions.

\subsection{Two Distinct Critiques}

\citet[pp.~344--346]{Goldberger1979} challenges the inference from a
variance decomposition to the effects and desirability of interventions.
Heritability describes variation under particular conditions; it does not,
by itself, establish the response to changing those conditions or the
benefits and costs of doing so. A small environmental variance component
can reflect little environmental variation rather than a small causal
response. Even if genetic effects operate entirely through biological
pathways, high heritability does not establish that those pathways cannot
be modified or their consequences compensated. Goldberger's eyeglasses
example illustrates why the genetic origin of a disadvantage does not
establish that remedying it is ineffective or prohibitively costly.

\citet[pp.~729--734]{Jencks1980} challenges the interpretation of genetic
and social causes as mutually exclusive. Genes can affect the environments
people select and the treatment they receive; those environments can then
affect outcomes. Thus a genetic variance component may include socially
mediated effects, and $1-h^2$ does not cap the causal importance of social
environments. His ``Genetic = Physical and Environmental = Social''
fallacy confuses the source of variation with the mechanism transmitting
its effects. Table~1 contrasts PKU (genetic/physical), lead poisoning
(nongenetic/physical), sexism (genetic differences eliciting social
treatment), and language spoken at home (nongenetic/social). Jencks also
qualifies these categories: lead exposure is socially patterned, while
PKU's consequences depend on diet, screening, and treatment. Entire
causal chains cannot be divided neatly into physical and social causes.

These are complementary emphases, not exclusive claims of intellectual
ownership: both authors discuss modifiability and policy. Neither critique
requires genetic estimates to be biased, and neither demonstrates that a
particular intervention works. The distinction is also developed in
\citet[secs.~3.4 and 3.6]{BenjaminEtAl2026}.

\subsection{``Direct'' Does Not Mean Narrowly Physiological}

Clark describes the direct path as one in which ``the genome dictates
child social aptitudes, independent of any interaction with the parents''
\citep[p.~286]{Clark2026}. This gives ``direct'' a meaning it does not
have in the statistical-genetic distinction between direct and indirect
genetic effects. As the potential-outcomes framework in
\citet[secs.~3.1 and 3.4]{BenjaminEtAl2026} makes explicit, the relevant
contrast changes the individual's own genotype while holding other
people's genotypes fixed; it does not hold fixed all environmental
responses to the individual. A child's genetically influenced behavior
may evoke parental attention, teacher responses, or peer treatment, which
then affect outcomes. Those pathways contribute to the effect of the
child's own genotype despite operating socially. An indirect parental
genetic effect instead originates in the parent's genotype and affects
the child other than through transmission of the child's alleles.

Clark therefore mistakenly treats direct genetic effects as necessarily
operating through the narrowly physiological pathways criticized by
Jencks. Evidence against parental genetic nurture, even if decisive,
would not establish that own-genotype effects bypass social mechanisms.
Nor would it rule out the other sources of confounding described in
Section~\ref{sec:confounding-beyond-nurture}. The two questions must be
kept separate: omitted-variable bias concerns identification of the
causal effect; social mediation concerns how that effect operates.

\subsection{Passages Vulnerable to the Critiques}

The following quotations identify the inferences at issue. Page numbers
refer to the supplied book draft \citep{Clark2026}; line wrapping is
normalized, but wording is retained. They are not, by themselves, evidence
that every empirical claim accompanying them is false.

\begin{enumerate}[label=\arabic*.]
\item ``To what extent are social outcomes such as occupations, income, wealth, longevity, and health socially versus genetically determined?'' (p.~5).

\emph{Jencks:} Genetic and social causation need not be competing explanations: social pathways can mediate genetic effects.

\item ``But social influences matter little. The orthodox belief in the primacy of social factors and the effectiveness of social interventions is wildly mistaken.'' (p.~6).

\emph{Both:} A genetic fit does not exclude social mediation or establish that interventions are ineffective.

\item ``The finding both in modern and historical England, and also in modern Denmark and Sweden, that genetic transmission can explain observed social outcomes, is indeed evidence of the meritocratic nature of these societies.'' (p.~7).

\emph{Jencks:} Discriminatory or otherwise nonmeritocratic social responses to genetically influenced characteristics can contribute to genetic effects.

\item ``This implies that at maximum genetics explains 60\% of adult outcomes, and other environmental factors 40\%.'' (p.~17).

\emph{Jencks:} Even granting the numerical decomposition, it does not divide causal importance into mutually exclusive genetic and environmental shares.

\item ``However, the environmental factors that matter to outcomes are not typically the observable shared environment twins experience -- the home environment, the education of parents, their parenting styles, school and neighborhood characteristics.'' (p.~17).

\emph{Both:} A small shared-environment component does not establish weak causal effects of named institutions or parenting; their effects may be genetically mediated or their observed variation limited.

\item ``The first is the direct path where the genome dictates child social aptitudes, independent of any interaction with the parents.'' (p.~286).

\emph{Jencks:} An effect of the child's own genotype can operate through parental responses to the child. It need not be independent of parent-child interaction, and is distinct from an effect of parental genotype holding the child's genotype fixed.

\item ``If direct additive genetic determination was the overwhelming cause of social outcomes, parental death at young ages for children should have little negative impact on child outcomes.'' (p.~189).

\emph{Both:} The genetic share of ordinary outcome variation does not by itself predict the effect of parental death; necessary or compensating rearing inputs and genotype-evoked responses require separate analysis.

\item ``The first implication of the importance of direct genetic transmission of social outcomes is that modern societies spend too much on education.'' (p.~9).

\emph{Goldberger:} Whether expenditure is excessive requires evidence on marginal benefits and costs, not a genetic variance share.

\item ``The second implication of direct genetic transmission is that social outcomes will change -- true educational attainment, social abilities -- only when the genetic endowment of the society changes.'' (p.~10).

\emph{Both:} A fixed genetic distribution does not fix outcome levels, environmental responses, or the effects of genes under different institutions.

\item ``But it does emphasize that if social causation is limited, the relative fiscal costs of different types of immigrants will endure across all future generations.'' (p.~291).

\emph{Goldberger, with Jencks relevant to the premise:} Limited social effects in the observed setting do not establish immutable relative fiscal costs under future institutions, tax rules or labor-market returns. The conditional wording must be retained.

\end{enumerate}

\subsection{Acknowledgments and Qualifications in Clark}

A fair assessment must also recognize the contrary passages. Some
explicitly state the relevant principle; others recognize an example or
alternative explanation without developing its general implications.
These are ten passages, not ten independent arguments or equally complete
acknowledgments of the critiques.

\begin{enumerate}[label=\arabic*.]
\item ``This does not in itself imply that social interventions cannot change social outcomes.'' (p.~80).

Preceding sentence supposes status mainly determined by genetic inheritance; explicitly rejects an automatic inference to intervention ineffectiveness.

\item ``The absence of convincing evidence for genetic nurture still leaves the issue of whether social interventions can be successful in changing outcomes for children in significant ways.'' (p.~86).

Especially strong: intervention remains an empirical question even on his assessment of genetic nurture.

\item ``The fact of genetic determination matters little, however, if social interventions can significantly change social outcomes.'' (p.~219).

Sets up his intervention-specific assessment of education; evaluate that evidence separately.

\item ``Thus how hereditable a trait is depends also on the range of variation in family environments. The smaller is this range of variation the more heritable will traits appear.'' (p.~14).

Acknowledges the distinction between observed environmental variation and the magnitude of environmental effects; retains source spelling hereditable.

\item ``Thus in poorer societies there are significant social class differences in height.'' (p.~14).

In context he attributes these differences to childhood access to food. Recognizes social determination of a physical outcome; not a complete statement of Jencks's critique.

\item ``Here the parental genotype creates a social environment for the child which favors better child social outcomes.'' (p.~80).

Describes genetic nurture as a possible socially mediated pathway, subsequently disputes its empirical importance. Does not cover all mediation through the child's own genotype.

\item ``Perhaps what is happening here is that there were sufficient social safety nets, even in nineteenth century England, that the loss of one or more parents was largely compensated for.'' (p.~191).

Allows an important environmental input to be replaced; small parental-loss effects need not mean care is causally unimportant.

\item ``Either nurture was not of great impact, or social mechanisms existed to replace lost parents.'' (p.~288).

Restates the compensation alternative in the conclusions, though he then argues persistence of parent-child correlations favors his interpretation.

\item ``Women were largely excluded from universities, and professional societies, before 1920.'' (p.~99).

Recognizes institutional restrictions on measured attainment; connects to Jencks's sexism example but is not an explicit discussion of mediated genetic effects.

\item ``Exceptions to this generalization, for example, would be the Jewish emigrees to England and the USA from Tsarist Russia, whose economic and social prospects were severely limited by social hostility and legal restrictions in their homeland.'' (p.~291, footnote~105).

Footnote 105 qualifies persistence of migrants' social rank; recognizes outcomes constrained by institutions and changed opportunities. Retains emigrees spelling.

\end{enumerate}

\subsection{The Inconsistency in Interpretation}

The strongest acknowledgments are on pp.~80, 86, and 219. It would be
inaccurate to say that Clark never recognizes the compatibility of genetic
inheritance with effective intervention, or that his entire case rests
on heritability alone. He discusses environmental pathways and presents
intervention-specific evidence. His height example recognizes the
dependence of heritability on environmental variation; his parental-loss
discussion recognizes that small observed effects can reflect replacement
of an important input. The examples involving women and Jewish migrants
also acknowledge institutional constraints on attainment.

The criticism is that these qualifications are not consistently carried
through to the claims that genetic transmission implies meritocracy,
excessive education spending, or genetic change as a necessary condition
for social improvement. Recognizing parental genetic nurture is narrower
than recognizing all social pathways through which an individual's own
genotype can affect outcomes. Likewise, rejecting that particular source
of confounding does not validate all the model's identifying restrictions.
The policy claims require evidence about particular interventions, their
benefits and costs, and the settings to which the findings generalize.
The genetic model cannot supply whatever that evidence leaves unresolved.
